% GENERATED FILE. Edit canonical lessons, then run npm run generate:books.
% canonical-source-hash: fnv1a64:9cbd20902a67d050
% canonical-lessons: JA-C257-nakushimashita, JA-C257-noriokureru, JA-C257-nihonjuu, JA-W257-juu-kanji, JA-R257-something-went-wrong

\chapter{Something Went Wrong}
\label{ch:ja-something-went-wrong}
\hlchaptermodality{257}

\begin{chapteropening}
\textbf{By the end of this chapter:} \emph{I can say what went wrong on a trip (saifu o nakushimashita; kasa ga kowaremashita; densha ni noriokuremashita), ask for help (tasukete kudasai), and say all over Japan (nihonjū).}

Complete JA-R257-something-went-wrong: report a lost wallet, a missed train and a broken umbrella, and ask for help.
\end{chapteropening}

\section[\texorpdfstring{\ja{なくしました} … (nakushimashita …)}{なくしました … (nakushimashita …)}]{\ja{なくしました} / \ja{こわれました} (nakushimashita / kowaremashita) — something went wrong}

\label{lesson:JA-C257-nakushimashita}



\begin{quote}
Say \textquotedblleft{}to lose\textquotedblright{} and \textquotedblleft{}to break down\textquotedblright{}. (\emph{nakusu}, \emph{kowareru}.) Then \textquotedblleft{}a wallet\textquotedblright{}. (\emph{saifu}.)
\end{quote}

\subsection*{You'll want to know: \ja{を} \ja{なくしました}, \ja{が} \ja{こわれました}}
\begin{quote}
\textbf{\ja{さいふを} \ja{なくしました。}}
\end{quote}

\begin{quote}
\emph{saifu o nakushimashita.} — \textquotedblleft{}I have lost my wallet.\textquotedblright{}
\end{quote}

\textbf{\ja{なくす}} is something you did, so the lost thing takes \textbf{\ja{を}}. \textbf{\ja{こわれる}} is something that happened by itself, so the broken thing takes \textbf{\ja{が}}:

\begin{quote}
\textbf{\ja{かさが} \ja{こわれました。}}
\end{quote}

\begin{quote}
\emph{kasa ga kowaremashita.} — \textquotedblleft{}My umbrella broke.\textquotedblright{}
\end{quote}

\begin{quote}
\textbf{\ja{きのう} \ja{かぎを} \ja{なくしました。}}
\end{quote}

\begin{quote}
\emph{kinō kagi o nakushimashita.} — \textquotedblleft{}I lost my key yesterday.\textquotedblright{}
\end{quote}

\begin{quote}
\textbf{\ja{でんわが} \ja{こわれました。}}
\end{quote}

\begin{quote}
\emph{denwa ga kowaremashita.} — \textquotedblleft{}The telephone stopped working.\textquotedblright{}
\end{quote}

\subsection*{You'll want to know: \ja{たすけて} \ja{ください}}
\textbf{\ja{たすけて} \ja{ください}} (\emph{tasukete kudasai}), the \textbf{-\ja{て}} form of \textbf{\ja{たすける}} with \textbf{\ja{ください}}, is \textquotedblleft{}please help me\textquotedblright{}:

\begin{quote}
\textbf{\ja{たすけて} \ja{ください。}}
\end{quote}

\begin{quote}
\emph{tasukete kudasai.} — \textquotedblleft{}Please help me.\textquotedblright{}
\end{quote}

\begin{quote}
\textbf{\ja{さいふを} \ja{なくしたので、こまっています。}}
\end{quote}

\begin{quote}
\emph{saifu o nakushita node, komatte imasu.} — \textquotedblleft{}I have lost my wallet, so I am in trouble.\textquotedblright{}
\end{quote}

\begin{quote}
\textbf{\ja{こうばんは} \ja{どこですか。}}
\end{quote}

\begin{quote}
\emph{kōban wa doko desu ka.} — \textquotedblleft{}Where is the police box?\textquotedblright{}
\end{quote}

\subsection*{Your turn}
Say these aloud:

\begin{itemize}
\raggedright
  \item \textquotedblleft{}I have lost my key\textquotedblright{} — \emph{kagi o nakushimashita}
  \item \textquotedblleft{}I have lost my ticket\textquotedblright{} — \emph{kippu o nakushimashita}
  \item \textquotedblleft{}I have lost my bag\textquotedblright{} — \emph{kaban o nakushimashita}
\end{itemize}

\begin{tcolorbox}[breakable,colback=teal!4,colframe=teal!35!black,title={Before you move on}]
\textquotedblleft{}My umbrella broke\textquotedblright{}? (\textbf{\ja{かさが} \ja{こわれました}}: it happened by itself, so \textbf{\ja{が}}.) \textquotedblleft{}I lost my key\textquotedblright{}? (\textbf{\ja{かぎを} \ja{なくしました}}.) \textquotedblleft{}Please help me\textquotedblright{}? (\textbf{\ja{たすけて} \ja{ください}}.)
\end{tcolorbox}
\section[\texorpdfstring{\ja{のりおくれる} (noriokureru)}{のりおくれる (noriokureru)}]{\ja{のりおくれる} (noriokureru) — to miss (a train or bus)}

\label{lesson:JA-C257-noriokureru}



\begin{quote}
Before the new word: say \textquotedblleft{}to ride\textquotedblright{} and \textquotedblleft{}to be late\textquotedblright{}. (\emph{noru}, \emph{okureru}.)
\end{quote}

\subsection*{You'll want to know: \ja{のりおくれる}}
\textbf{\ja{のりおくれる}} — \emph{noriokureru} — \textquotedblleft{}to miss\textquotedblright{} a train or a bus: to be too late to get on. Polite \textbf{\ja{のりおくれます}}, \emph{noriokuremasu}; an ichidan verb, like \textbf{\ja{おくれる}}. Past (polite): \textbf{\ja{でんしゃに} \ja{のりおくれました}}, \emph{densha ni noriokuremashita}.

\begin{quote}
\textbf{\ja{でんしゃに} \ja{のりおくれました。}}
\end{quote}

\begin{quote}
\emph{densha ni noriokuremashita.} — \textquotedblleft{}I missed the train.\textquotedblright{}
\end{quote}

\begin{quote}
\textbf{\ja{あさねぼうを} \ja{しました。}}
\end{quote}

\begin{quote}
\emph{asanebō o shimashita.} — \textquotedblleft{}I overslept.\textquotedblright{}
\end{quote}

\begin{quote}
\textbf{\ja{それで、でんしゃに} \ja{のりおくれました。}}
\end{quote}

\begin{quote}
\emph{sore de, densha ni noriokuremashita.} — \textquotedblleft{}And so I missed the train.\textquotedblright{}
\end{quote}

\begin{cousinweb}
\textbf{\ja{のりおくれる}} (\emph{noriokureru}) is a native Japanese word (wago), inherited rather than borrowed from Chinese: the stem of \textbf{\ja{のる}} (\emph{noru}), \textquotedblleft{}to ride\textquotedblright{}, and \textbf{\ja{おくれる}} (\emph{okureru}), \textquotedblleft{}to be late\textquotedblright{}. JMdict glosses it \textquotedblleft{}to miss (a train, bus, etc.)\textquotedblright{} (JMdict).
\end{cousinweb}

\subsection*{Your turn}
Say these aloud:

\begin{itemize}
\raggedright
  \item \emph{noriokureru}
  \item \textquotedblleft{}I missed the train\textquotedblright{} — \emph{densha ni noriokuremashita}
  \item \textquotedblleft{}I missed the subway\textquotedblright{} — \emph{chikatetsu ni noriokuremashita}
\end{itemize}

\begin{tcolorbox}[breakable,colback=teal!4,colframe=teal!35!black,title={Before you move on}]
\textquotedblleft{}To miss a train\textquotedblright{}? (\textbf{\ja{のりおくれる}}.) Which particle goes with the train? (\textbf{\ja{に}}.)
\end{tcolorbox}
\section[\texorpdfstring{\ja{日本中} (nihonjū)}{日本中 (nihonjū)}]{\ja{日本中} (nihonjū) — all over Japan}

\label{lesson:JA-C257-nihonjuu}



\begin{quote}
Before the new word: say \textquotedblleft{}a Japanese person\textquotedblright{} and \textquotedblleft{}all year round\textquotedblright{}. (\emph{nihonjin}, \emph{nenjū}.)
\end{quote}

\subsection*{You'll want to know: \ja{日本中}}
\textbf{\ja{日本中}} — \emph{nihonjū} — \textquotedblleft{}all over Japan\textquotedblright{}, the whole country through. The \textbf{\ja{じゅう}} at its end is the one in \textbf{\ja{年}\ja{じゅう}}.

\begin{quote}
\textbf{\ja{日本中}\ja{に} \ja{ともだちが} \ja{います。}}
\end{quote}

\begin{quote}
\emph{nihonjū ni tomodachi ga imasu.} — \textquotedblleft{}I have friends all over Japan.\textquotedblright{}
\end{quote}

You write \textbf{\ja{日本}}. The third sign is new, read \textbf{\ja{じゅう}} (\emph{jū}) here; the next lesson writes it.

\begin{cousinweb}
\textbf{\ja{日本中}} (\emph{nihonjū}) is Sino-Japanese (kango): three Chinese characters read with their on'yomi, readings borrowed from Chinese. They are \emph{ni}, \textquotedblleft{}sun\textquotedblright{} (Middle Chinese \emph{nyit}), \emph{hon}, \textquotedblleft{}origin\textquotedblright{} (\emph{pwonX}) and \emph{jū}, \textquotedblleft{}in, inside, middle\textquotedblright{} (\emph{trjuwng}), the voiced \emph{chū} of \textbf{\ja{年}\ja{じゅう}} (KANJIDIC2; Middle Chinese in Baxter's transcription). JMdict glosses it \textquotedblleft{}throughout Japan\textquotedblright{} (JMdict).
\end{cousinweb}

\subsection*{Your turn}
Say these aloud:

\begin{itemize}
\raggedright
  \item \emph{nihonjū}
  \item \textquotedblleft{}I have friends all over Japan\textquotedblright{} — \emph{nihonjū ni tomodachi ga imasu}
  \item \textquotedblleft{}There are shrines all over Japan\textquotedblright{} — \emph{nihonjū ni jinja ga arimasu}
\end{itemize}

\begin{tcolorbox}[breakable,colback=teal!4,colframe=teal!35!black,title={Before you move on}]
\textquotedblleft{}All over Japan\textquotedblright{}? (\textbf{\ja{日本中}}, \emph{nihonjū}.) How is the new sign read here? (\textbf{\ja{じゅう}}, \emph{jū}.)
\end{tcolorbox}
\section[\texorpdfstring{\ja{中} (jū)}{中 (jū)}]{\ja{中} — the middle, four strokes}

\label{lesson:JA-W257-juu-kanji}



\begin{quote}
Say \emph{nihonjū}, all over Japan. Then call up \textbf{\ja{口}}, a mouth: three strokes, a box.
\end{quote}

\subsection*{Script — the shape}
\begin{quote}
\textbf{\ja{中}} — \emph{jū}
\end{quote}

Its parts:

\begin{itemize}
\raggedright
  \item \ja{口} (box)
  \item a long centre stem through it
\end{itemize}

Four strokes, three lifts:

\begin{enumerate}
\raggedright
  \item draw the left side of the box from the top down
  \item lift, draw the top from left to right, then without lifting turn down the right side
  \item lift, close the base of the box from left to right
  \item lift, draw the long centre stem from above the box straight down through it
\end{enumerate}

The box is \textbf{\ja{口}}, wider than it is tall, and the stem comes last, running through it from above to below. Here it is read \textbf{\ja{じゅう}} (\emph{jū}), as in \textbf{\ja{日本中}}; this lesson asks only for that reading.

\subsection*{Your turn}
\begin{enumerate}
\raggedright
  \item Find \textbf{\ja{中}} in \textbf{\ja{日本中}} and say \emph{nihonjū}.
  \item Copy \textbf{\ja{中}} three times, saying \emph{jū} at each one.
  \item Find it once more, in \textbf{\ja{日本中}\ja{に} \ja{ともだちが} \ja{います}} (\emph{nihonjū ni tomodachi ga imasu}), and say the line.
\end{enumerate}

\begin{tcolorbox}[breakable,colback=teal!4,colframe=teal!35!black,title={Before you move on}]
Which sign do you know inside it? (\textbf{\ja{口}}.)

Source: \href{https://github.com/KanjiVG/kanjivg/blob/master/kanji/04e2d.svg}{KanjiVG stroke-order data}: KanjiVG, kanji/04e2d.svg, kanji \ja{中} (U+4E2D), 4 stroke paths kvg:04e2d-s1 to s4 (Ulrich Apel and contributors, CC BY-SA 3.0). Only the order and direction are taken from it.
\end{tcolorbox}
\section[(review)]{Review — something went wrong}

\label{lesson:JA-R257-something-went-wrong}



\begin{quote}
Say \textquotedblleft{}I have lost my wallet\textquotedblright{}. (\emph{saifu o nakushimashita}.) Then \textquotedblleft{}I missed the train\textquotedblright{}. (\emph{densha ni noriokuremashita}.)
\end{quote}

\subsection*{Your turn}
Say these aloud:

\begin{itemize}
\raggedright
  \item My umbrella broke.
  \item I lost my key yesterday.
  \item Where is the police box?
  \item I missed the subway.
  \item There are shrines all over Japan.
  \item What time does the next train leave?
  \item Please turn right at the bank.
\end{itemize}

\begin{tcolorbox}[breakable,colback=teal!4,colframe=teal!35!black,title={Before you move on}]
Lost: \textbf{\ja{を}} or \textbf{\ja{が}}? (\textbf{\ja{を}}: \textbf{\ja{さいふを} \ja{なくしました}}.) Broke? (\textbf{\ja{が}}: \textbf{\ja{かさが} \ja{こわれました}}.) \textquotedblleft{}Please help me\textquotedblright{}? (\textbf{\ja{たすけて} \ja{ください}}.) \textquotedblleft{}All over Japan\textquotedblright{}? (\textbf{\ja{日本中}}.)
\end{tcolorbox}
